\documentclass[10pt,a4paper]{amsart}
\usepackage[margin=19mm]{geometry}
\usepackage{amsmath,amssymb,mathtools}
\usepackage[scr=rsfs,cal=euler]{mathalfa}
\usepackage{xfrac}
\usepackage[unicode,hidelinks,bookmarks=false]{hyperref}
\newcommand{\ie}{i.e.\ }
\newcommand{\cf}{cf.\ }
\newcommand{\resp}{respectively\ }
\newcommand{\Subsection}{\subsection}
\providecommand{\nobreakdash}{\nobreak\hbox{-}}
\setlength{\emergencystretch}{2em}
\multlinegap0pt
\usepackage{bm}
\usepackage{bbm}
\usepackage{dsfont}
\usepackage{xspace}
\usepackage{cancel}
\usepackage{mathtools}
\usepackage{amsthm}
\makeatletter
\@ifundefined{thm}{\newtheorem{thm}{Theorem}}{}
\@ifundefined{claim}{\newtheorem{claim}[thm]{Claim}}{}
\makeatother
\newcommand{\R}{\mathbb{R}}
\newcommand{\osc}{\mathfrak{osc}}
\newcommand{\sign}[1]{\mathrm{sign}(#1)}
\title[Convergence rate of $\ell^p$-energy minimization on graphs: ]{Convergence rate of $\ell^p$-energy minimization on graphs: sharp polynomial bounds and a phase transition at $p=3$}
\author{Gideon Amir, Fedor Nazarov,, Yuval Peres}
\date{Source: arXiv:2508.19411v1 (CC BY 4.0)}
\begin{document}
\maketitle

\noindent\emph{Source and licence.} Convergence rate of $\ell^p$-energy minimization on graphs: sharp polynomial bounds and a phase transition at $p=3$. Version arXiv:2508.19411v1 (CC BY 4.0), distributed under the Creative Commons Attribution 4.0 International licence, as recorded in the arXiv licence metadata for this exact version and shown on the abstract page https://arxiv.org/abs/2508.19411v1. Full rights notice: licenses/arxiv-2508.19411-CC-BY-4.0.txt.\par\smallskip

\noindent\textit{Editorial scope.} Counted unit: the Thm whose source label is \texttt{None} in the article (source0.tex lines 2122--2124, source label \texttt{None}), delivered together with the article's own complete original proof of it (source0.tex lines 2126--2274, one \texttt{proof} environment of 149 lines). No mathematics is written, repaired or re-derived: statement and proof are the article's own text line for line. Required context retained uncounted: eq:psuper (lines 717--719). Every definition, display and label of those lines is retained unchanged. Omitted from this package and not redistributed: 1--716, 720--2121, 2125--2125, 2275--2723. Nothing inside a retained excerpt is omitted or altered: every retained body is delivered exactly as the article's lines are written, so no omission receipt of any kind accompanies this package. Citations: the retained excerpts carry no citation command at all, so no bibliography entry of the article is transcribed and no reference is redistributed. Reference adaptation: none was needed and none was made; every reference inside the delivered unit resolves inside it, so no unresolved reference or citation is produced. Printed slips kept literal: no printed word, symbol, hypothesis or number of the counted statement or of its proof was altered. Layout and class: the article's own class is \texttt{amsart} with packages including fancyhdr, booktabs, amssymb, latexsym, amsthm, enumitem, color, amsmath, array, mathtools, bbm, nicefrac, upquote, tikz; the wrapper is the lane's pinned \texttt{amsart} editorial wrapper at 10pt on A4, which restates the theorem-like environments the retained text needs and adds the article's own macro definitions that the retained text needs (\texttt{\textbackslash R}, \texttt{\textbackslash osc}, \texttt{\textbackslash sign}), with extra wrapper packages (bm, bbm, dsfont, xspace, cancel, mathtools). The delivered numbering is the unit's own. Assets: no retained span contains a figure environment, a graphics-inclusion command, a PDF-page inclusion, a table or a TikZ picture; no image file is copied into this package, the package declares no asset dependency and no image file is redistributed with it. Licence: version arXiv:2508.19411v1 of this article is distributed under the Creative Commons Attribution 4.0 International licence, as recorded in the arXiv licence metadata for this exact version and shown on the abstract page https://arxiv.org/abs/2508.19411v1. A full rights notice is delivered at licenses/arxiv-2508.19411-CC-BY-4.0.txt. Characters: every retained excerpt is pure ASCII, so the delivered engine, XeLaTeX with its default Latin Modern text font, reports no missing character.
\begin{equation} \label{eq:psuper}
0 \le \Psi_f'(f(v))=p\sum_{w \sim v}|f(v)-f(w)|^{p-1}\sign(f(v)-f(w)) \,.
\end{equation}

    \begin{thm}
        Run the $\ell^p$-energy minimizing dynamics on the graph $H_{d,n}$ with initial opinions $f_0$ defined as follows:  $f_0=0$  on all vertices $w_{i,j}$ and on $[-n,0)$, at the center $f_0(0)=1/2$,  and     $f_0=1$ on  all $u_{i,j}$ and on $(0,n]$.  Then for every large enough $n$ and $d$, every updating sequence, and every $t<\frac{1}{25e} 2^\frac{-p}{p-1}n^\frac{2p-1}{p-1}d^\frac{1}{p-1}$,  we have $\osc(f_t)\geq \frac{1}{2}$.
    \end{thm}

\begin{proof}
Let $\delta =   e\left(\frac{2}{nd}\right)^{\frac{p}{p-1}} $
(this choice of $\delta$ will emerge from the proof). 
Define a new initial profile $h_0:V\rightarrow [0,1]$ by 
    \begin{eqnarray*}
        &&h_0(w_{i,j}) = (j-1) \delta \quad 1\leq i\leq m\, , 1\leq j \leq d ,\\
        &&h_0(-n)=e^{-1} d^\frac{p}{p-1}\delta=  (2/n)^{\frac{p}{p-1}}\, ,\\
        &&h_0(i) = h_0(-n) + (i+n)\frac{1-h_0(-n)}{n} \quad -n<i\leq 0 \, , \\
        &&h_0(i) = 1 \quad 0<i\leq n \, ,\\
        &&h_0(u_{i,j}) = 1 \quad 1\leq i\leq m\, , 1\leq j \leq d .\\
    \end{eqnarray*}
    
Fix an updating sequence. As in the previous example, we let $h_t$ denote the evolution of $h_0$ under the $\ell^p$-energy minimizing dynamics with the same update sequence. 
Since $h_0\geq f_0$, it is enough to prove that 

\begin{equation}\label{eq:large D main}  \forall t<\frac{1}{25e} 2^\frac{-p}{p-1}n^\frac{2p-1}{p-1}d^\frac{1}{p-1},\quad  \forall i\le m, \quad  \min_j h_t(w_{i,j})<\frac{1}{4} \, ,\end{equation}  
as this would imply that $\osc(f_t)\geq \frac{1}{2}$ for such $t$.

We will need  some further notation. 
Given  a function $h:V\rightarrow \R$, we denote by $\widetilde{h}$ the function obtained by re-ordering the values of $h$ on each clique $W_i$, so that the sequence $\{\widetilde{h}(w_{i,j})\}_{j=1}^d$ is nondecreasing in $j$ for every $1\leq i\leq m$. 
We denote by $k_i(t)$ the number of updates in $W_i$ by time $t$. 

We will deduce (\ref{eq:large D main}) from the following 
claim:
\begin{claim}
          For every integer $t\in \big[0, \frac{m}{6e}d^{\frac{p}{p-1}}\big]$, we have
    \begin{itemize}
        \item [(a)]
        For all $0\leq j\leq n$, 
\begin{equation}\label{eq:star}
            h_t(-j)\leq h_0(-j)  \, ;
        \end{equation}
        \item [(b)] For all $1\leq i\leq m$ and all $1\leq j\leq d$,
\begin{equation}\label{eq:starstar}
\widetilde{h_t}(w_{i,j})\leq \min\bigl(h_0(-n), \, h_0(w_{i,j})+k_i(t)\delta \bigr) \, .
\end{equation}        
    \end{itemize}
\end{claim}
\begin{proof}
    We will prove both clauses together by induction on $t$.  For $t=0$ the claim is trivial, so we move to the induction step and assume that $t+1 \le  \frac{m}{6e}d^{\frac{p}{p-1}}$ and the inequalities \eqref{eq:star} and \eqref{eq:starstar} hold at time $t$. 
    %Using symmetry, we may assume that $h_t=\widetilde{h_t}$.
    %is enough to show the induction step when we replace   $h_t(w_{i,j})$ by $h_0(w_{i,j})+k_i(t)\delta$ for all $1\leq i\leq m$ and $1\leq j\leq d$. 
    
    We start by showing that \eqref{eq:starstar} holds for time $t+1$. The inequality
    \begin{equation} \label{clear1}
     \forall i,j \quad  \; h_{t+1}(w_{i,j}) \le h_0(-n)  
    \end{equation}
    is immediate from the induction hypothesis and monotonicity of the dynamics. 
    
    We may assume that the vertex $v_{t+1}$ updated at time $t+1$ is in  $W_i$ for some $i \le m$, since otherwise \eqref{eq:starstar} clearly continues to hold.
    %moreover, we may assume that $v_{t+1}$  is the vertex $w_{i,1}$ where $h_t$ attains its minimum on $W_i$ (since the resulting values of $\widetilde{h_{t+1}}$ dominate pointwise the values obtained if we update another vertex in the same clique). 
    Reorder  $W_i=\{\zeta_j\}_{j=1}^d$ so that $j \mapsto h_{t}(\zeta_j)$ is nondecreasing, i.e., for all $j \in \{1,\ldots,d\}$, we have
    \begin{equation} \label{enumi0}
    h_{t}(\zeta_j)=\widetilde{h_{t}}(w_{i,j}) \,. 
    \end{equation} 
We also reorder  $W_i=\{z_j\}_{j=1}^d$ so that $j \mapsto h_{t+1}(z_j)$ is nondecreasing and  we have
    \begin{equation} \label{enumi}
    \forall j \in \{1,\ldots,d\}, \quad \;h_{t+1}(z_j)=\widetilde{h_{t+1}}(w_{i,j}) \,. 
    \end{equation} 
    
    Suppose that $v_{t+1}=z_\ell \in W_i$. For   $j\in  \{1,\ldots,\ell-1\}$
    we have 
    $z_j \in \{\zeta_j,  \zeta_{j+1}\}$,  so 
    by the induction hypothesis and \eqref{enumi0},
    $$h_{t+1}(z_j)  \le h_t(\zeta_{j+1})\le h_0(w_{i,j+1})+k_i(t)\delta=h_0(w_{i,j})+k_i(t+1)\delta \,.$$
    For   $j\in  \{\ell+1,\ldots,d\} $
    we have $z_j\in \{\zeta_j,  \zeta_{j-1}\}$,
    whence 
    $$h_{t+1}(z_j)  \le h_t( \zeta_j)\le h_0(w_{i,j})+k_i(t+1)\delta\,.$$
     In view of \eqref{clear1} and \eqref{enumi}, this verifies   \eqref{eq:starstar}  at time $t+1$ for all $j \ne \ell$. 
    
    Similarly, if $j=\ell<d$, then
    $z_{\ell+1}\in \{\zeta_\ell,  \zeta_{\ell+1}\}$, so
    $$h_{t+1}(z_\ell) \le h_{t+1}(z_{\ell+1}) \le h_{t}(\zeta_{\ell+1}) \le h_0(w_{i,\ell+1})+k_i(t)\delta = h_0(w_{i,\ell})+k_i(t+1)\delta \,,$$ 
    verifying \eqref{eq:starstar}  at time $t+1$ for   $j = \ell$ if $\ell<d$.

The only remaining case is $j=\ell=d$, and for this case we must check that
\begin{equation} \label{eq:l-equals-d} h_{t+1}(z_d) \stackrel{?}{\leq} h_0(w_{i,d})+k_i(t+1)\delta \,.
\end{equation}
When verifying this, we may assume that $h_0(w_{i,d})+k_i(t+1)\delta < h_0(-n)$, since otherwise   \eqref{clear1} implies that \eqref{eq:l-equals-d} holds. 
    By the superharmonicity criterion \eqref{eq:psuper}, it is enough to prove that
    \begin{eqnarray*}  
\sum_{j=2}^{d}\bigl(h_0(w_{i,d})+k_i(t+1)\delta -\widetilde{h_t}(w_{i,j})\bigr)^{p-1} 
  \stackrel{?}{\geq} \bigl(h_t(-n)-h_0(w_{i,d})-k_i(t+1)\delta \bigr)^{p-1} .
    \end{eqnarray*}
Recalling the induction hypotheses
\eqref{eq:star} and \eqref{eq:starstar}, as well as   the definition of $h_0$, the preceding inequality   would follow  if we show that
\begin{equation}
    \sum_{j=2}^{d}\bigl((d-j+1)\delta \bigr)^{p-1} \stackrel{?}{\geq}  h_0(-n)^{p-1}\, .
\end{equation}
    Comparing the sum to an integral, it suffices to verify that
    \begin{equation}
        \frac{(d-1)^p}{p}\delta^{p-1}\stackrel{?}{\geq} \Bigl(e^{-1}d^\frac{p}{p-1}\delta \Bigr)^{p-1} \, .
    \end{equation}
    It is therefore enough to prove that
    \begin{equation}
        \frac{(d-1)^p}{p}\stackrel{?}{\geq}  {e^{1-p}}d^p \, ,
    \end{equation}
    which is equivalent to 
    \begin{equation}
        \left(\frac{d-1}{d}\right)^p\stackrel{?}{\geq}pe^{1-p} \, .
    \end{equation}
    The last inequality holds for all $d>d(p)$, since   $pe^{1-p}<1$ for $p>1$.

    We now move to show that \eqref{eq:star} holds at time $t+1$. This only has to be checked if at time $t+1$ we update one of the vertices $-n,\ldots, 0$, and only at the updated vertex. For any $j\in \{-n+1,\dots,0\}$, this holds by monotonicity of the dynamics and the definition of $h_0$, so we are left with the only interesting case of updating vertex $-n$. Thus we need to show that $h_{t+1}(-n)\leq h_0(-n)$. For this to hold it is enough to check that 
   \begin{equation} 
       \sum_{i=1}^m \sum_{j=1}^d \Bigl(h_0(-n)  - \widetilde{h_t}(w_{i,j})\Bigr)^{p-1} 
 \stackrel{?}{\geq} \big(h_t(1-n)-h_0(-n)\big)^{p-1} \,.
 %\left(\frac{1}{n}\right)^{p-1} \, .
   \end{equation}
    By the induction hypothesis,  it suffices to prove that 
\begin{equation}\label{eq: 3}
        \sum_{i=1}^n \sum_{j=1}^d \max \bigl(0 \,,   h_0(-n)-h_0(w_{i,j})-k_i(t) \delta \bigr)^{p-1}\stackrel{?}{\geq} n^{1-p} \,.
    \end{equation}
    Let $S_t= \{i\leq m \, : \, k_i(t)\leq \frac{2t}{m} \}$.
    For  $d>d(p)$,  we have
    $d+\frac{1}{3e}d^{\frac{p}{p-1}} <\frac{1}{2e} d^{\frac{p}{p-1}}$, so if also $t< \frac{m}{6e}d^{\frac{p}{p-1}}$ and $i \in S_t$, then  
    $$ h_0(w_{i,j})+k_i(t) \delta \le h_0(w_{i,j})+\frac{2t}{m} \delta \le \frac{1}{2} h_0(-n) \,.$$
     Thus to verify \eqref{eq: 3}, it is enough to check   that
    \begin{equation}
        \sum_{i\in S_t} d\left( \frac{1}{2} h_0(-n)\right)^{p-1}\stackrel{?}{\geq} n^{1-p} \, .
    \end{equation}
    
     Since $\sum_{1\leq i\leq m}k_i(t)\leq t$, we have $|S_t|\geq \frac{m}{2}$, so it suffices to show that
    \begin{equation}
        \frac{md}{2} \Bigl(\frac12 h_0(-n) \Bigr)^{p-1}\stackrel{?}{\geq} n^{1-p} \,. 
    \end{equation}
    Recalling that $md=n$ and that $h_0(-n)=\frac{1}{e} d^\frac{p}{p-1}\delta$, the last inequality will follow for large $d$ if we show that
    \begin{equation}
        \frac{n}{2}\left(\frac{1}{2e}d^\frac{p}{p-1}\delta\right)^{p-1}\stackrel{?}{\geq} n^{1-p} \, .
    \end{equation}
    This, indeed, holds for  
    $\delta =   e\left(\frac{2}{nd}\right)^{\frac{p}{p-1}}$.   The proof of the claim is complete. \end{proof}
Let $t_*:= \Big\lfloor\frac{m}{6e}d^{\frac{p}{p-1}} \Big\rfloor \,.$
The claim implies that $\widetilde{h_{t_*}} \leq h_0 + h_0(-n)$ (pointwise). Another induction gives that for all integers $\ell \ge 1$, we have
\[ \widetilde{h_{\ell t_*}}\leq h_0+\ell h_0(-n) \,.\]  

We deduce that 
$$  \forall t<\left\lfloor\frac{1}{4h_0(-n)} \right\rfloor t_* \, \,,\quad  \forall i\le m \,,\quad \; \min_j h_t(w_{i,j})<\frac{1}{4} \, .$$ Since $$\left\lfloor\frac{1}{4h_0(-n)} \right\rfloor t_* \geq \frac{1}{25e} 2^\frac{-p}{p-1}n^\frac{2p-1}{p-1}d^\frac{1}{p-1} \, $$ for $d>d(p)$, the theorem is proved.

% If $\ell t_* \le t <(\ell+1)t_*$, then the preceding inequality yields that
% $$\forall i \le 2n, \quad h_{kn}(w_i) \le k n^\frac{-p}{p-1} \,.$$
% Since at most $n$ updates take place in the time interval $[kn,t]$, 
% the inequality \eqref{eq:h2 grows slow} follows, completing the proof.

    
%     To deduce \eqref{eq:large D main} from the above claim, we simply note that $k_i^*(t) = \frac{2t}{m}$ for at least half the clusters, and that in any such cluster we have $h_t(w_{i,1})\leq \frac{2t}{m}\cdot \delta$. Thus $h_t(w_{i,1})\leq \frac{1}{4}$ for all $t\leq \frac{m}{8\delta} = \frac{1}{24e}n^\frac{2p-1}{p-1}d^\frac{1}{p-1}$. 
% This completes the proof of the Theorem. 
\end{proof}

\end{document}
